\documentclass[10pt,a4paper]{amsart}
\usepackage[margin=19mm]{geometry}
\usepackage{amsmath,amssymb,mathtools}
\usepackage[scr=rsfs,cal=euler]{mathalfa}
\usepackage{xfrac}
\usepackage[unicode,hidelinks,bookmarks=false]{hyperref}
\newcommand{\ie}{i.e.\ }
\newcommand{\cf}{cf.\ }
\newcommand{\resp}{respectively\ }
\newcommand{\Subsection}{\subsection}
\providecommand{\nobreakdash}{\nobreak\hbox{-}}
\setlength{\emergencystretch}{2em}
\multlinegap0pt
\usepackage{amssymb}
\usepackage{mathtools}
\usepackage[shortlabels]{enumitem}
\usepackage{xcolor}
\usepackage{algorithm}\usepackage{algpseudocode}
\usepackage{float}
\usepackage{bm}
\newtheorem{definition}{Definition}
\newtheorem{remark}{Remark}
\newtheorem{lemma}{Lemma}
\newtheorem{theorem}{Theorem}
\newcommand*{\bbB}{\mathbb{B}}
\newcommand*{\bbD}{\mathbb{D}}
\newcommand*{\calC}{\mathcal{C}}
\newcommand*{\calD}{\mathcal{D}}
\newcommand*{\calE}{\mathcal{E}}
\newcommand*{\calF}{\mathcal{F}}
\newcommand*{\calS}{\mathcal{S}}
\newcommand*{\calU}{\mathcal{U}}
\newcommand*{\calV}{\mathcal{V}}
\DeclareMathOperator{\fl}{fl}
\DeclareMathOperator{\gr}{gr}
\title{Hilbert's 16th problem for arrangements of curves on a surface}
\author{Giacomo Maletto}
\date{Source: arXiv:2606.21449v1 (CC BY 4.0)}
\begin{document}
\maketitle

\noindent\emph{Source and licence.} Hilbert's 16th problem for arrangements of curves on a surface. Version arXiv:2606.21449v1 (CC BY 4.0), distributed by arXiv under the Creative Commons Attribution 4.0 International licence, http://creativecommons.org/licenses/by/4.0/, as recorded in the arXiv licence metadata for this exact version and shown on the abstract page https://arxiv.org/abs/2606.21449v1. Full licence notice: \texttt{licenses/arxiv-2606.21449-CC-BY-4.0.txt}.\par\smallskip

\noindent\textit{Editorial scope.} Counted unit: the article's theorem thm:main (source0.tex lines 773--779). The article's complete original proof of it is delivered with it (source0.tex lines 781--820, one \texttt{proof} environment of 40 lines). No mathematics is written, repaired or re-derived: the statement and the proof are the article's own lines, in the article's own notation, and no retained line is substituted or re-flowed.  Retained uncounted as the source-local context the proof needs: lines 732--746; lines 754--756; lines 758--760; lines 766--768.  Nothing was omitted from any retained span, so the package carries no omission record.  No retained line contains a citation, so the wrapper carries no bibliography and no bibliography entry.  No retained line contains a figure environment, an image-inclusion command or drawing code, so no image is delivered and the package declares no asset dependency.  Editorial formatting only, in addition: an AMS \texttt{amsart} preamble that restates the article's own declarations and macros used by the retained lines, namely the environments \texttt{definition}, \texttt{remark}, \texttt{lemma}, \texttt{theorem} on separate counters, together with the standard packages those lines need.  The retained environments print in this document as Definition 1, Remark 1, Lemma 1, Theorem 1 in order of appearance, whereas the article prints the same items with the numbering of its own full document.  The complete native source of the article as distributed in the arXiv e-print for this exact version is bundled unchanged as \texttt{sources/203270/source0.tex}.
\begin{definition}\label{def:combcurve}
    Let $(V,E,F)$ be a combinatorial cell complex on $\calS$. A \textbf{combinatorial curve} transverse to $(V,E,F)$ is given by the following data $(n,W,T)$:
    \begin{itemize}
        \item for every edge $e\in E$, a non-negative integer $n(e)$ such that $n(f):=\sum_{i=1}^{l(f)}n(e_i)$ is even for all faces $f\in F$;
        \item for every face $f\in F$, a Dyck word $W_f$ of length $n(f)$;
        \item for every $g\in\sum_{f\in F}[n(f)/2+1]$, a rooted tree $T_g$, called the \textbf{floating tree} on $g$.
    \end{itemize}

    From a curve $\calD$ transverse to $(\calV,\calE,\calF)$ we construct the following combinatorial curve $(n,W,T)$ transverse to $(V,E,F)$:
    \begin{itemize}
        \item for every edge $e\in\calE$, let $n(e):=\#(\calD\cap e)$ be the number of points of intersection of $e$ with the curve $\calD$;
        \item for every face $f\in\calF$, let $W_f=W(\varphi_f^{-1}(\calD))$ be the Dyck word associated to the curve $\varphi_f^{-1}(\calD)$ on $\bbD^2$;
        \item for every face $f\in\calF$, consider the ground faces $g\in\calF_{\gr(\calD)},g\subseteq f$ in the order in which each $\overline{\varphi_f^{-1}(g)}$ is first met by traveling along $\partial\varphi_{\bbD^2}$. If $g$ is the $i$-th, let $T_{(f,i)}$ be the rooted tree formed by the regions of $\calD$ contained inside $g$, where the root is the unique ground region whose closure contains $\partial g$.
    \end{itemize}  
\end{definition}

\begin{remark}\label{rem:act}
    Let $\calD$ be a curve intersecting $(\calV,\calE,\calF)$ transversely, and consider the induced combinatorial curve $(n,W,T)$. Given a homeomorphism $\eta$, the curve $\eta(\calD)$ will have an induced combinatorial curve $(n',W',T')$ which depends uniquely on $(n,W,T)$ and the functions $\eta_V$, $\eta_E$ and $\eta_F$. In this way, the group of automorphisms of $(V,E,F)$ acts on the combinatorial curves transverse to $(V,E,F)$. Clearly, a homeomorphism $\eta$ fixes $(\calV,\calE,\calF)$ if and only if it induces the trivial action.
\end{remark}

\begin{lemma}\label{lem:disk}
    Up to homeomorphism fixing $\partial\bbD^2$ pointwise, a curve $\calC$ on $\bbD^2$ whose every component meets the boundary $\partial\bbD^2$ is determined by the set $\calC\cap\partial\bbD^2$ and the Dyck word $W(\calC)$.
\end{lemma}

\begin{lemma}\label{lem:float}
    Up to homeomorphism fixing $\partial\bbD^2$ pointwise, a curve $\calC$ on $\bbD^2$ disjoint from $\partial\bbD^2$ is determined by the rooted tree formed by the connected components of $\bbB^2\setminus\calC$, where the root is the unique component whose closure contains $\partial\bbD^2$.
\end{lemma}

\begin{theorem}\label{thm:main}
    The construction in Definition \ref{def:combcurve} defines a bijection
    \begin{equation}\label{eq:main}
        \frac{\{\calD\mid\calD\pitchfork(\calV,\calE,\calF)\}}{\text{homeomorphism that fixes $(\calV,\calE,\calF)$}}\to\{(n,W,T)\mid(n,W,T)\pitchfork(V,E,F)\}
    \end{equation}
    between curves $\calD$ transverse to $(\calV,\calE,\calF)$ up to homeomorphism that fixes $(\calV,\calE,\calF)$, and combinatorial curves transverse to $(V,E,F)$.
\end{theorem}

\begin{proof}
    The function is well-defined: suppose $\calD$ and $\calD'$ are equivalent through a homeomorphism $\eta$ that fixes $(\calV,\calE,\calF)$. Then, by Remark \ref{rem:act}, they have the same associated combinatorial curve.

    To prove surjectivity, choose $n_e$ points on each edge $e$. For each face $f$, draw the unique non-crossing arc system prescribed by $W_f$. Since the endpoint labels agree on both sides of every edge, these arc systems glue to a curve $\gr(\calD)$. Finally, in each ground face $g$ insert floating ovals according to the rooted tree $T_g$.

    We now focus on injectivity. Suppose $\calD,\calD'$ are two curves meeting $(\calV,\calE,\calF)$ transversely to which we associate the same data $(n,W,T)$. Then up to a homeomorphism that fixes $(\calV,\calE,\calF)$ we can assume $\calD\cap e=\calD'\cap e$ for all $e\in\calE$. To do this, fix an edge $e$. Our goal is to find an open set $\calU\subset\calS$ such that
    \begin{itemize}
        \item $\calU\cap e$ is a connected set containing both $\calD\cap e,\calD'\cap e$ and such that the following equalities hold:
        \begin{equation}\label{eq:closureU}
            \overline{\calU}\cap\calV=\varnothing,\quad\overline{\calU}\cap\bigcup\calE=\overline{\calU}\cap e;
        \end{equation}
        \item it is ``easy'' to find a homeomorphism $\overline{\calU}\simeq\overline{\calU}$ fixing the boundary $\partial\calU$ pointwise, $\calU\cap e$ setwise and that sends $\calD\cap e$ into $\calD'\cap e$.
    \end{itemize}
    This homeomorphism fixes $(\calV,\calE,\calF)$ and achieves what we want. We construct $\calU$ explicitly, distinguishing two cases:
    \begin{itemize}
        \item[$\underline{e\not\subset\partial\calS}$:] by definition of 2-manifold, there exist exactly two tuples $(f,i),(f',i')$ such that $\partial_i(f),\partial_{i'}(f')=\pm e$. Recall that
        \[\partial\varphi_f|_{(t_{i-1},t_i)}\colon(t_{i-1},t_i)\to e\]
        is a homeomorphism. Pick values $a,b\in(t_{i-1},t_i)$ such that \[(\partial\varphi_f|_{(t_{i-1},t_i)})^{-1}((\calD\cap e)\cup(\calD'\cap e))\subset(a,b).\]
        Define
        \[T=\varphi_f(\{(r\cos(2\pi t),r\sin(2\pi t))\mid r\in(0,1),t\in(a,b)\})\subset\Delta_i(f).\]
        Next, let
        \[a'=(\partial\varphi_{f'}|_{(t_{i'-1},t_{i'})})^{-1}(\partial\varphi_f|_{(t_{i-1},t_i)}(a)),\quad b'=(\partial\varphi_{f'}|_{(t_{i'-1},t_{i'})})^{-1}(\partial\varphi_f|_{(t_{i-1},t_i)}(b))\]
        and likewise define
        \[T'=\varphi_{f'}(\{(r\cos(2\pi t),r\sin(2\pi t))\mid r\in(0,1),t\in(a',b')\})\subset\Delta_{i'}(f').\]
        Finally, let
        \[\calU=T\cup T'\cup\partial\varphi_f((a,b)).\]
        Then, we have that $\calU$ is open, $\calU\cap e=\partial\varphi_f((a,b))$ is connected containing both $\calD\cap e,\calD'\cap e$, which satisfies the equalities \eqref{eq:closureU}. Moreover,
        \[(\overline{\calU},\calU,\calU\cap e)\simeq(\bbD^2,\bbB^2,0\times(-1,1))\]
        so it is easy to find a homeomorphism $\overline{\calU}\to\overline{\calU}$ with the desired properties.
        \item[$\underline{e\subset\partial\calS}$:] there exists a unique tuple $(f,i)$ with $\partial_i(f)=\pm e$. Pick $a,b$ and define $T$ as before. Let
        \[\calU=T\cup\partial\varphi_f((a,b)).\]
        Now we instead have that
        \[(\overline{\calU},\calU,\calU\cap e)\simeq(\bbD^2\cap\{y\ge0\},\bbB^2\cap\{y\ge0\},0\times[0,1)),\]
        but everything works as before.
    \end{itemize}

    Next we show that for each $f\in\calF$, there is a homeomorphism $\eta_f\colon\overline{f}\to\overline{f}$ fixing $\partial f$ such that $\eta_f(\gr(\calD)\cap\overline{f})=\gr(\calD')\cap\overline{f}$. Indeed, the curves $\varphi_f^{-1}(\gr(\calD)\cap\overline{f}), \varphi_f^{-1}(\gr(\calD')\cap\overline{f})$ on $\bbD^2$ have the same Dyck word $W_f$ associated to them in Lemma \ref{lem:disk}, so there is a homeomorphism $\eta\colon\bbD^2\to\bbD^2$ fixing $\partial\bbD^2$ sending the first curve to the second. Since $\eta$ fixes $\partial\bbD^2$ pointwise, it preserves the fibers of $\varphi_f$. Therefore $\eta$ descends to a homeomorphism $\eta_f\colon\overline{f}\to\overline{f}$ satisfying $\eta_f\circ\varphi_f=\varphi_f\circ\eta$, which has the desired properties. Since each $\eta_f$ is the identity on $\partial f$, these glue together to form a homeomorphism $\eta\colon\calS\to\calS$ fixing $(\calV,\calE,\calF)$ and such that $\eta(\gr(\calD))=\gr(\calD')$.

    We can now suppose $\gr(\calD)=\gr(\calD')$. Similarly to before, using Lemma \ref{lem:float} we conclude that for each ground face $g\in\calF_{\gr(\calD)}$, the floating part $\fl(\calD)\cap g$ is uniquely specified by the rooted tree $T_g$ up to homeomorphism of $\overline{g}$ fixing its boundary.
\end{proof}

\end{document}
